\section{Introduction}\label{section::Introduction}
% \textcolor{blue}{
% Ideas: Laws of random probabilities are now everywhere: in BNP, to model data-sets, to model distributions of images, in functional data analysis, etc. In this work: two-sample test.}
% \textcolor{blue}{
% Related work: consider if there are available hypothesis testing in arbitrary metric spaces and come up with convincing motivation of why our approach is more natural. Mention that one could embed the probabilities in an Hilbert space through the quantile function but a) this works only on $\R$ and b) find (standard) motivation for preferring working on the space itself; e.g. look for motivation of paper using the standard wasserstein for testing.}

The two-sample test is a fundamental statistical procedure used to test the null hypothesis that two datasets are generated from the same underlying distribution. Specifically, given two sets of samples, the goal is to assess whether their associated distributions are statistically indistinguishable. The hypotheses are formulated as follows:
\begin{itemize}
    \item Null hypothesis ($H_0$): The two distributions are identical.
    \item Alternative hypothesis ($H_1$): The two distributions are distinct.
\end{itemize}
Classic examples include the Kolmogorov–Smirnov test, which compares empirical cumulative distribution functions (CDFs) derived from independent and identically distributed (i.i.d.) samples. Other well-known univariate methods include the classical $t$-test and the Wilcoxon rank-sum test. As modern applications frequently involve multivariate observations in $\mathbb{R}^d$, several extensions have been developed, including kernel-based methods \cite{gretton2012kernel} and approaches based on the Wasserstein distance \cite{ramdas2017wasserstein, wang2021two}.

Increasingly, however, data objects have become more complex, often appearing as probability measures themselves. Settings where observed samples are modeled as probability measures are now prevalent \cite{petersen2022modeling}; examples include mortality rates across different countries, income distributions within various regions, and distributional summaries of functional connectivity in the brain. It is important to mention that, in the majority of cases, these probability measures are latent and are instead estimated from the raw observations they generate. Two-sample testing, in the case of probability measures as samples, naturally motivates the study of the laws of random probability measures, that is, probability measures defined on the space of probability measures.

This problem was previously addressed in \cite{dubey2019frechet} using Fr\'echet means and variances. However, that approach has limitations: specifically, if the laws are not uniquely characterized by their Fr\'echet means and variances, the test may fail to provide a high rejection rate when the null hypothesis is false. In contrast, we adopt a fully nonparametric perspective, making no distributional assumptions on the underlying laws of the RPMs. Furthermore, theoretical guarantees are provided directly in terms of the raw observations used to estimate the underlying probability measures.

In this work, we adopt a distance-based approach. This framework first requires a metric to quantify the distance between two datasets. The core intuition is that if the observed distance is sufficiently "large," it is "unlikely" that the underlying distributions are the same. To formalize the terms "large" and "unlikely," one must determine a threshold based on a pre-specified significance level $\theta$. This threshold is defined such that the probability of the observed distance exceeding it under the null hypothesis is at most $\theta$. Consequently, the testing procedure is to reject the null hypothesis if the observed distance exceeds this threshold. 



% While extensions of two-sample tests to general metric spaces already exist (see Section 3.7 in \cite{van1996weak}), they are not directly applicable to our setting, since we must also account for the sampling scheme. In particular, given two laws of random probability measures $Q^1, Q^2$, we do not assume i.i.d.\ samples of probability measures from $Q^1$ and $Q^2$. Instead, we observe i.i.d.\ samples of probability measures from other laws that are close to the original ones $Q^1$ and $Q^2$ (see details in Section \ref{sec::hier_sampling}). Formally, the extended setting compared to the one in \cite{van1996weak} is described as follows: Let $P,Q$ be probability measures on a general metric space $\mathcal{X}$. Consider sequences of probability measures $(P_m)_{m \geq 1}$ and $(Q_m)_{m \geq 1}$ such that  
% \[
% P_m \overset{w}{\to} P \quad \text{and} \quad Q_m \overset{w}{\to} Q.
% \]  
% Instead of observing  
% \[
% X_1,\ldots, X_n \sim P \quad \text{and} \quad Y_1,\ldots, Y_n \sim Q,
% \]  
% we observe  
% \[
% X^m_1,\ldots, X^m_n \overset{\mathrm{i.i.d.}}{\sim} P_m 
% \quad \text{and} \quad 
% Y^m_1,\ldots, Y^m_n \overset{\mathrm{i.i.d.}}{\sim} Q_m,
% \]  
% for some fixed $m \geq 1$.


Our objective is to establish theoretical guarantees for both the false positive rate (the probability of incorrectly rejecting the null hypothesis when it is true) and the power (the probability of correctly rejecting the null hypothesis when it is false). 

The rest of the paper is organized as follows. In Section \ref{sec::hier_sampling}, we formalize the sampling scheme and discuss its connections to exchangeability. Section \ref{sec::distances} reviews existing distance functions on laws of random probability measures. In Section \ref{sec::rademacher}, we develop a theoretical threshold based on Rademacher complexity and illustrate its practical limitations through simulations. As a remedy, in Section \ref{section::permutation}, we introduce a practical threshold based on a permutation approach and establish its theoretical guarantees. Finally, in Section \ref{section::simulations}, we benchmark our method against existing approaches in the literature on simulated datasets and apply it to a mortality dataset.
